\section{A free root and a planar perturbation}
\label{fr:section}\label{sec:free-root}

We now study an edge norm defined by bounded scalar potentials whose
root value is free. The root-inclusive sibling model also permits such a
root value, but records it as a root coefficient; here it contributes no
coordinate. The rootless model $K_{\rm s}$ instead fixes the root potential
at zero. We give a direct construction of witnesses in the edge coordinates: inward jumps detect
the first edges beyond a finite head, and a perturbation in the Euclidean
plane detects the remaining edges.

Let $T=\mathbb N^{<\omega}$, with root $o$, and let
$T^+=T\setminus\{o\}$.  Identify $s\in T^+$ with the edge from its
parent $s^-$ to $s$, and write $\operatorname{ch}(v)$ for the children of
$v$.  Define
\begin{equation}
 \mathcal B=\left\{b:T\longrightarrow[-1,1]:
 \sum_{w\in\operatorname{ch}(v)}|b_w-b_v|^2\le1
 \quad(v\in T)\right\}.
 \label{fr:tests}
\end{equation}
For $b\in\mathcal B$ and $x\in c_{00}(T^+)$ put
\begin{equation}
 f_b(x)=\sum_{s\in T^+}x_s(b_s-b_{s^-}),
 \qquad \|x\|_{\mathrm e}=\sup_{b\in\mathcal B}f_b(x),
 \label{fr:norm}
\end{equation}
and let $X_{\mathrm e}$ be the completion.  In particular, there is no
dummy edge above $o$ and no condition $b_o=0$.

Every increment in \eqref{fr:tests} has absolute value at most one.
For each $s\in T^+$, the potential equal to one on the subtree rooted
at $s$ and zero elsewhere has just one nonzero increment, at $s$.
Together with symmetry of $\mathcal B$, these observations give
\[
 |x_s|\le\|x\|_{\mathrm e}\le\sum_{u\in T^+}|x_u|,
 \qquad \|e_s\|_{\mathrm e}=1.
\]
Thus \eqref{fr:norm} is a norm and all coordinate functionals are
continuous.  Each $f_b$ extends to $X_{\mathrm e}$ with norm at most
one, and the same supremum formula remains valid there: both sides are
$1$-Lipschitz and agree on the dense space $c_{00}(T^+)$.

Call $D\subset T$ a finite initial subtree if it is finite, contains
$o$, and contains every predecessor of each of its nodes.  Given a
potential $b$, retain its labels on $D$ and freeze them at the last node
in $D$ along every path leaving $D$.  This preserves admissibility and
sets precisely the increments outside $D$ to zero.  Consequently the
coordinate projection
\[
 P_Dx=\sum_{s\in D\cap T^+}x_se_s
\]
is contractive, and $Q_D=I-P_D$ has norm at most two.  The space
\begin{equation}
 F_D=\ker P_D
   =\overline{\operatorname{span}}\{e_s:s\notin D\}
 \label{fr:tail}
\end{equation}
is closed and has finite codimension.  For the second equality, apply
$Q_D$ to finite approximants of a vector in $\ker P_D$.
Finally, a finite vector supported in $D$ has a norming potential frozen
outside $D$.  Indeed the possible labels on $D$ form a compact set:
they lie in $[-1,1]^D$ and satisfy the finitely many closed child
constraints within $D$; freezing gives an admissible extension of every
such choice.

\begin{theorem}
\label{thm:free-root-modulus}
The real Banach space $X_{\mathrm e}$ defined by \eqref{fr:tests} and
\eqref{fr:norm} satisfies
\begin{equation}
 \widehat\delta_{\|\cdot\|_{\mathrm e}}(t)
 \ge\frac{\sqrt{1+t^2/576}-1}{4}\qquad(0<t<1).
 \label{fr:modulus}
\end{equation}
It admits no equivalent AUC norm.  More precisely, if
$\alpha\|\cdot\|_{\mathrm e}\le N\le\beta\|\cdot\|_{\mathrm e}$,
where $0<\alpha\le\beta<\infty$, then
$\overline\delta_N(\alpha/(4\beta))=0$.
\end{theorem}

The midpoint estimate is proved first for finite centers and tails.
The head norming potential already gives a lower bound for both
endpoints.  We therefore need only increase one endpoint by an amount
depending on $t$.  The first-exit and deeper-edge cases supply this
increase by different extensions of that potential.

\begin{proof}
Fix $0<t<1$.  Let $x$ be a finite unit vector supported in an initial
subtree $D$, and choose a norming potential $b^0$ frozen outside $D$.
For every $y\in F_D$,
\begin{equation}
 f_{b^0}(x)=1,\qquad f_{b^0}(y)=0,
 \qquad \|x+ty\|_{\mathrm e},\ \|x-ty\|_{\mathrm e}\ge1.
 \label{fr:two-lower-bounds}
\end{equation}
Suppose first that $y\in F_D$ is finite and $\|y\|_{\mathrm e}=1$.
Let $E=\{r\notin D:r^-\in D\}$ be the first-exit nodes.  Split
$y=y_E+y_L$ into its coordinates on $E$ and on the deeper edges.

\emph{First-exit case.}
Suppose some $c\in\mathcal B$ satisfies $|f_c(y_E)|\ge1/2$.
For each $v\in D$ set
\[
 L_v=\left(\sum_{r\in\operatorname{ch}(v)\setminus D}|y_r|^2\right)^{1/2}.
\]
Cauchy--Schwarz in each sibling group gives $\sum_{v\in D}L_v\ge1/2$.
Choose the inward sign $\sigma_v=-1$ when $b^0_v\ge0$ and
$\sigma_v=1$ when $b^0_v<0$.  For $\varepsilon\in\{-1,1\}$ define
\[
 L_{v,\varepsilon}=
 \left(\sum_{\substack{r\in\operatorname{ch}(v)\setminus D\\
                     \varepsilon\sigma_v y_r>0}}|y_r|^2\right)^{1/2}.
\]
Since $L_{v,1}+L_{v,-1}\ge L_v$, one common sign $\varepsilon$ has
$\sum_vL_{v,\varepsilon}\ge1/4$.

For $0<h<1$, multiply all head labels by $1-h^2$.  At a selected exit
$r$ from $v$, assign the jump
\[
 d_r=\frac{\varepsilon h y_r}{L_{v,\varepsilon}};
\]
assign zero at the other exits, omit groups with
$L_{v,\varepsilon}=0$, and freeze each exit label throughout its
descendant subtree.  Every nonzero jump has sign $\sigma_v$ and
magnitude at most $h$.  It points toward zero from the head label;
if it crosses zero, its resulting magnitude is at most $h<1$.
Thus all labels remain in $[-1,1]$.  At a parent in $D$ the internal
and exit increments have total square sum at most
\[
 (1-h^2)^2+h^2\le1,
\]
and outside $D$ all further increments vanish.  This is an admissible
potential.  Its evaluations on $x$ and $y$ are respectively $1-h^2$
and $\varepsilon h\sum_vL_{v,\varepsilon}$.  Hence
\begin{equation}
 \max\{\|x+ty\|_{\mathrm e},\|x-ty\|_{\mathrm e}\}
 \ge1-h^2+th/4.
 \label{fr:exit}
\end{equation}
Taking $h=t/8$ gives the lower bound $1+t^2/64$.

\emph{Deeper-edge case.}
If the preceding case fails, every admissible test has
$|f_c(y_E)|<1/2$.  Choose $c\in\mathcal B$ with $f_c(y)>5/6$;
then $f_c(y_L)>1/3$.  We shall preserve the norming head on the
horizontal axis and use the vertical coordinate to detect this deeper
contribution.

Consider labels $B_v\in\mathbb R^2$ satisfying
\begin{equation}
 |B_v|_2\le1,\qquad
 \sum_{w\in\operatorname{ch}(v)}|B_w-B_v|_2^2\le1.
 \label{fr:plane-admissibility}
\end{equation}
Every scalar projection of $B$ onto a unit vector belongs to
$\mathcal B$.  Therefore the linear map
\[
 A_Bz=\sum_{s\in T^+}z_s(B_s-B_{s^-})
\]
satisfies $|A_Bz|_2\le\|z\|_{\mathrm e}$ for finite $z$, and extends
contractively to $X_{\mathrm e}$.

Set $B_v=(b^0_v,0)$ on $D$.  For each first-exit node $r$, put
$u=b^0_{r^-}$ and, throughout the subtree rooted at $r$, define
\begin{equation}
 k=\frac18,\qquad z_v=k(c_v-c_r),\qquad
 B_v=\bigl(u\sqrt{1-z_v^2},z_v\bigr).
 \label{fr:plane-extension}
\end{equation}
The first-exit increment is zero because $z_r=0$.
Also $|z_v|\le2k=1/4$, and
$|B_v|_2^2=u^2(1-z_v^2)+z_v^2\le1$.
For fixed $u$ and $c_r$, the map
\[
 H(a)=\bigl(u\sqrt{1-k^2(a-c_r)^2},k(a-c_r)\bigr),
 \qquad -1\le a\le1,
\]
has
\begin{equation}
 |H'(a)|_2^2
 =k^2+\frac{u^2k^4(a-c_r)^2}{1-k^2(a-c_r)^2}
 \le\frac{k^2}{1-4k^2}=\frac1{60}.
 \label{fr:plane-lipschitz}
\end{equation}
It is consequently $1$-Lipschitz.  At a parent outside $D$ this same
map applies to all its children, so the child-increment constraint
follows from that for $c$.  At a parent in $D$, the original head
increments remain admissible and all first-exit increments are zero.
This verifies \eqref{fr:plane-admissibility} at the boundary as well
as in every descendant subtree.

Now $A_Bx=(1,0)$, whereas the second coordinate of $A_By$ is
$k f_c(y_L)>1/24$.  The first exits contribute zero to that coordinate.
The parallelogram identity gives
\begin{align}
 \max\{\|x+ty\|_{\mathrm e}^2,\|x-ty\|_{\mathrm e}^2\}
 &\ge\frac{|A_Bx+tA_By|_2^2+|A_Bx-tA_By|_2^2}{2}\\
 &\ge1+t^2/576.
 \label{fr:deep}
\end{align}

\emph{Completion and arbitrary centers.}
Put $G(t)=\sqrt{1+t^2/576}-1$.  Since
$G(t)\le t^2/1152<t^2/64$, both cases give a maximum endpoint norm
at least $1+G(t)$.  Combining this with \eqref{fr:two-lower-bounds}
yields an average at least $1+G(t)/2$.
Finite unit vectors are dense in the unit sphere of $F_D$, by
\eqref{fr:tail} followed by normalization, so continuity gives this
same estimate for every unit $y\in F_D$.
For a unit $u\in F_D$, the function
\[
 r\longmapsto\frac{\|x+ru\|_{\mathrm e}+\|x-ru\|_{\mathrm e}}2
\]
is convex and even, hence nondecreasing on $[0,\infty)$.  The estimate
therefore also holds for every $y\in F_D$ with $\|y\|_{\mathrm e}\ge1$.
Given an arbitrary unit center, choose a finite unit center within
$G(t)/4$ and use its subspace $F_D$.  The reverse triangle inequality
loses at most $G(t)/4$ in the average.  Taking the infimum over these
tail directions, the supremum over finite-codimensional subspaces, and
then the infimum over unit centers proves \eqref{fr:modulus}, with the
direction convention of \eqref{eq:averaged-modulus}.

\emph{The renorming obstruction.}
For nonroot $s$, let $p_s$ be the sum of the edges on the path from
$o$ to $s$.  Telescoping gives $f_b(p_s)=b_s-b_o$, and the singleton
increment test gives the lower bound in
\begin{equation}
 1\le\|p_s\|_{\mathrm e}\le2.
 \label{fr:paths}
\end{equation}
For every fixed parent $s$ and finite scalar family $(a_n)$,
\begin{equation}
 \left\|\sum_{n=1}^m a_ne_{s^\frown n}\right\|_{\mathrm e}
 \le\left(\sum_{n=1}^m|a_n|^2\right)^{1/2}.
 \label{fr:siblings}
\end{equation}
This follows directly from the sibling constraint in
\eqref{fr:tests}.  If $\varphi\in X_{\mathrm e}^*$, applying
\eqref{fr:siblings} with $a_n=\varphi(e_{s^\frown n})$ shows
\[
 \sum_{n=1}^m|\varphi(e_{s^\frown n})|^2\le\|\varphi\|^2;
\]
thus $e_{s^\frown n}\to0$ weakly against the whole continuous dual.
These increments have norm one, and
$p_{s^\frown n}=p_s+e_{s^\frown n}$.
Start at any fixed nonroot $v$ and use the tree
$(p_{v^\frown s})_{s\in T}$.  By \eqref{fr:paths} and
\eqref{fr:siblings}, Lemma~\ref{lem:bounded-tree} applies with
$m=1$, $M=2$ and $c=1$.  Its obstruction scale is exactly
$\alpha c/(2\beta M)=\alpha/(4\beta)$, proving the final assertion.
\end{proof}
